% !TEX program = xelatex
\documentclass{extarticle}
\usepackage{amsmath}
\usepackage{graphicx}
\graphicspath{{media/}}
\usepackage{fontspec}
\usepackage{caption}
\captionsetup{font=small,labelfont=bf}
\setmainfont[
    Ligatures=TeX,
    UprightFont=DavidCLM-Medium.otf,
    BoldFont=DavidCLM-Bold.otf,
    ItalicFont=DavidCLM-MediumItalic.otf,
    BoldItalicFont=DavidCLM-BoldItalic.otf
]{David CLM}
\newfontfamily\hebrewfont[
    Script=Hebrew,
    Ligatures=TeX,
    UprightFont=DavidCLM-Medium.otf,
    BoldFont=DavidCLM-Bold.otf,
    ItalicFont=DavidCLM-MediumItalic.otf,
    BoldItalicFont=DavidCLM-BoldItalic.otf
]{David CLM}
\newfontfamily\hebrewfonttt[
    Script=Hebrew,
    UprightFont=MiriamMonoCLM-Book.ttf,
    BoldFont=MiriamMonoCLM-Bold.ttf,
    ItalicFont=MiriamMonoCLM-BookOblique.ttf,
    BoldItalicFont=MiriamMonoCLM-BoldOblique.ttf
]{Miriam Mono CLM}
\renewcommand{\familydefault}{\rmdefault}
\setlength{\parindent}{0bp}
\setlength{\parskip}{0.35em}
\usepackage{float}
\usepackage[pdfusetitle,
 bookmarks=true,bookmarksnumbered=false,bookmarksopen=false,
 breaklinks=false,pdfborder={0 0 1},backref=false,colorlinks=false]
 {hyperref}
\usepackage{siunitx}
\sisetup{mode=math, per-mode=fraction, bracket-unit-denominator=true, uncertainty-mode=separate, compound-units=single}
\DeclareSIUnit{\microfarad}{\ensuremath{\mu\mathrm{F}}}

\makeatletter
%%%%%%%%%%%%%%%%%%%%%%%%%%%%%% User specified LaTeX commands.
\usepackage[a4paper,margin=2cm,footskip=0.25in]{geometry}
\renewcommand{\labelitemi}{$\bullet$}
\usepackage{amssymb}
\usepackage{booktabs}
\usepackage{array}
\usepackage{tikz}
\usetikzlibrary{arrows.meta}
\renewcommand{\arraystretch}{1.18}
\makeatother

\usepackage{polyglossia}
\setdefaultlanguage{hebrew}
\setotherlanguage{english}
\begin{document}
\title{\texorpdfstring{פרומגנטיות וטמפרטורת קירי \\[0.4em] \Large מחברת מעבדה}{פרומגנטיות וטמפרטורת קירי — מחברת מעבדה}}
\author{איתי פלדמן \and טום בלכר}
\date{5 במאי 2026}

\maketitle

\tableofcontents
\bigskip

\part{שבוע 1}
\label{part:week1}

\section{מטרת הניסוי}
\begin{itemize}
\item מדידת הפרמאביליות של פרומגנט (ברזל).
\item מדידת פרמאביליות הריק.
\end{itemize}

\section{רקע תיאורטי}

\subsection{\texorpdfstring{השדות $B,H,M$ ופרמאביליות}{השדות B,H,M ופרמאביליות}}
בחומר שבו $\vec{M}=\chi\vec{H}$ ($\chi$ הסוספטיביליות המגנטית),
\begin{equation}
\vec{B}=\mu_{0}(\vec{H}+\vec{M})=\mu_{0}(1+\chi)\vec{H}\equiv\mu\vec{H},
\end{equation}
ומכאן
\begin{equation}
\mu_{\mathrm{r}}\equiv\frac{\mu}{\mu_{0}}=\frac{B}{\mu_{0}H}.
\end{equation}
בחומר לינארי $\mu_{\mathrm{r}}$ קבוע; בפרומגנט $\mu_{\mathrm{r}}=\mu_{\mathrm{r}}(H)$ ויש לולאת חשל.

\subsection{עקומת המגנוט ולולאת חשל}
אינטראקציית השחלוף מפצלת את החומר למתחמים מגנטיים המיושרים פנימית, שסכומם הווקטורי מתאפס ללא שדה. עם הגברת $H$ מתרחבים המתחמים שבכיוון השדה עד הרוויה (איור~\ref{fig:hysteresis}).

\begin{figure}[H]
    \centering
    \begin{minipage}[t]{0.49\linewidth}
        \centering
        \includegraphics[width=\linewidth]{ferromagnetic_domains.pdf}
        \caption{סכמה של מתחמים מגנטיים: ללא שדה חיצוני המתחמים מכוונים בכיוונים שונים, ואילו בשדה חזק מתקבלת רוויה מגנטית בכיוון השדה.}
        \label{fig:ferromagnetic-domains}
    \end{minipage}\hfill
    \begin{minipage}[t]{0.49\linewidth}
        \centering
        \includegraphics[width=\linewidth]{hysteresis.jpg}
        \caption{לולאת חשל סכמטית של חומר פרומגנטי, להגדרת רמננטיות, קוארסיביות ורוויה. הנתונים שלנו כוללים רק מעטפת שיא־לשיא ולא לולאה מלאה.}
        \label{fig:hysteresis}
    \end{minipage}
\end{figure}

\subsection{חילוץ פרמאביליות הריק}
עבור טבעת ברזל באורך מגנטי $L$ עם רווח בעובי $L'$ ממולא נחושת ($\mu_{\mathrm{Cu}}\approx\mu_{0}$), חוק אמפר נותן
\begin{equation}
\oint\vec{H}\cdot\mathrm{d}\vec{\ell}=H_{\mathrm{Fe}}\,L+H_{\mathrm{gap}}\,L'=NI.
\end{equation}
בהנחת רציפות $B_{\mathrm{Fe}}=B_{\mathrm{gap}}\equiv B$ (רלוונטית ברוויה),
\begin{equation}\label{eq:linear}
\frac{NI}{B}=\frac{1}{\mu_{0}}\,L'+\frac{L}{\mu_{\mathrm{Fe}}}.
\end{equation}
אם $B$ נשמר קבוע בסריקת $L'$, השיפוע הוא $1/\mu_0$.

\section{מערכת הניסוי}
טבעת ברזל עם שני סלילים (איור~\ref{fig:circuit}). הסליל הראשוני מוזן משנאי משתנה דרך נגד $R_x$, כך ש-$V_x\propto I_x\propto H$. הסליל המשני מחובר לאינטגרטור $RC$ ($R_y,C$), כך ש-$V_y\propto B$. המתחים $V_x,V_y$ מוזנים לצירי $X,Y$ של אוסילוסקופ המציג את הלולאה ישירות.

\subsection{ציוד}
\begin{itemize}
\item טבעת ברזל עם שני סלילים, $N_1=N_2=250$.
\item סט לוחיות נחושת.
\item שנאי משתנה, נגד $R_x$ בטור עם הסליל הראשוני, אמפרמטר לבקרת מגבלת זרם, ואינטגרטור $RC$.
\item אוסילוסקופ \textenglish{Agilent InfiniiVision DSO7012A} ומולטימטר \textenglish{HP 34401A}.
\end{itemize}

תרשים המערכת נראה באיור 
\ref{fig:circuit};
קבועי הכיול והגאומטריה מופיעים בטבלה
\ref{tab:apparatus-week1} (שגיאות לפי \ref{tab:uncertaintiesA}).

\begin{center}
\begin{minipage}[c]{0.44\linewidth}
\centering
\scriptsize
\setlength{\tabcolsep}{3pt}
\begin{tabular}{@{}rcc@{}}
\toprule
גודל & סימון & ערך \\
\midrule
\multicolumn{3}{@{}r@{}}{\textbf{נתוני תדריך}} \\
מספר ליפופים בכל סליל & $N\equiv N_1=N_2$ & $250$ \\
נגד האינטגרטור & $R_y$ & $\SI{11.10(3)}{\kilo\ohm}$ \\
קבל האינטגרטור & $C$ & $\SI{20.100(2)}{\microfarad}$ \\
\addlinespace[0.25em]
\multicolumn{3}{@{}r@{}}{\textbf{מדידות בניסוי}} \\
גובה הסליל & $h_{\mathrm{coil}}$ & $\SI{8.0}{\centi\meter}$ \\
גובה הטבעת & $h_{\mathrm{ring}}$ & $\SI{9.0}{\centi\meter}$ \\
היקף לולאת אמפר & $L$ & $\SI{0.4800(8)}{\meter}$ \\
נגד מדידת הזרם & $R_x$ & $\SI{2.9990(43)}{\ohm}$ \\
שטח חתך אפקטיבי & $A$ & $\SI{16.0(2)}{\centi\meter\squared}$ \\
\bottomrule
\end{tabular}
\captionof{table}{קבועי שבוע~1. גובה הסליל, גובה הטבעת, ההיקף והחתך נמדדו בסרגל.}
\label{tab:apparatus-week1}
\end{minipage}\hfill
\begin{minipage}[c]{0.53\linewidth}
\centering
\includegraphics[width=0.88\linewidth]{experiment.png}
\captionof{figure}{מערכת הניסוי: הסליל הראשוני נותן
$V_{x}\propto H$,
והסליל המשני מחובר לאינטגרטור כך ש-%
$V_{y}\propto B$.}
\label{fig:circuit}
\end{minipage}
\end{center}

\subsection{\texorpdfstring{גזירת $H(V_{x})$ ו-$B(V_{y})$}{גזירת H(Vx) ו-B(Vy)}}
חוק אמפר $HL=N_1 I_x$ ביחד עם $I_x=V_x/R_x$ נותן
\begin{equation}\label{eq:Hcal}
H=\frac{N_{1}}{L\,R_{x}}\,V_{x}.
\end{equation}
חוק פאראדיי $V_{\mathrm{ind}}=-N_2 A\,dB/dt$ ובחירת $\omega R_yC\gg1$ הופכים את הקבל לאינטגרטור,
\begin{equation}\label{eq:Bcal}
V_{y}\approx\frac{1}{R_{y}C}\int V_{\mathrm{ind}}\,\mathrm{d}t=\frac{N_{2}A}{R_{y}C}\,B\;\Longrightarrow\;B=\frac{R_{y}C}{N_{2}A}\,V_{y}.
\end{equation}
בפועל $\tau=R_yC=\SI{0.223}{\second}$ ו-$f=\SI{50}{\hertz}$, כך ש-$\omega\tau\approx70$.

הקריאות $\Delta V_x,\Delta V_y$ נמדדות שיא־לשיא, ולכן בחזרה על~\eqref{eq:Hcal}--\eqref{eq:Bcal} עם $\Delta V$ נכנס גורם $1/2$:
\begin{equation}\label{eq:HBpp}
H=\frac{N_{1}}{2\,L\,R_{x}}\,\Delta V_{x},\qquad B=\frac{R_{y}C}{2\,N_{2}A}\,\Delta V_{y}.
\end{equation}

\section{אי־ודאויות}
לכל גודל מדוד $x$,
$\sigma_x=\sqrt{(\sigma_{x,\mathrm{res}}+\sigma_{x,\mathrm{instr}}^2+\sigma_{x,\mathrm{other}}^2}$ (איברים לא רלוונטיים מושמטים). שגיאות המערכת מופיעים בטבלה~\ref{tab:uncertaintiesA}. 
$R_x$
 נמדד בתחום $\SI{100}{\ohm}$
  של המולטימטר ונניח כי
  $R_y$
  נמדד בתחום 
  $\SI{10}{\kilo\ohm}$
  (לא נתונה לו שגיאה אך נניח שנמדד עם אותו מולטימטר בטווח המתאים); השגיאה בקיבול נתונה.

\begin{table}[H]
\centering
\scriptsize
\setlength{\tabcolsep}{2pt}
\begin{tabular}{@{}cccccc@{}}
\toprule
גודל & מכשיר & $q$ (רזולוציה) & שגיאת מכשיר & שגיאה אחרת & שגיאה כוללת \\
\midrule
$a_L,b_L,a_A,b_A$ & סרגל & $\SI{1}{\milli\meter}$ & --- & --- & $q/\sqrt{12}$ \\
$L'$ & קליבר & $\SI{0.05}{\milli\meter}$ & --- & --- & $q/\sqrt{12}$ \\
$R_x$ & מולטימטר & $\SI{0.001}{\ohm}$ & $0.010\%\,R+0.004\%\cdot\SI{100}{\ohm}$ & --- & $\sqrt{(q/\sqrt{12})^2+\sigma_{\mathrm{instr}}^2}$ \\
$R_y$ & מולטימטר & $\SI{0.1}{\kilo\ohm}$ & $0.010\%\,R+0.001\%\cdot\SI{10}{\kilo\ohm}$ & --- & $\sqrt{(q/\sqrt{12})^2+\sigma_{\mathrm{instr}}^2}$ \\
$C$ & קבל & --- & --- & $0.010\%\,C$ & $\sigma_{\mathrm{other}}$ \\
$\Delta V_x,\Delta V_y$ & אוסצילוסקופ & $\SI{0.5}{\milli\volt}$ & $2.4\%\,|\Delta V|$ & --- & $\sqrt{(q/\sqrt{12})^2+\sigma_{\mathrm{instr}}^2}$ \\
\bottomrule
\end{tabular}
\caption{שגיאות המערכת: שגיאות המכשירים נלקחו מהתדריכים המתאימים להם בפרמטרים הלוונטיים לניסוי.}
\label{tab:uncertaintiesA}
\end{table}

הגדלים הנגזרים $L=2(a_L+b_L)$ ו-$A=a_Ab_A$ מורכבים משתי קריאות סרגל בלתי תלויות. אי־ודאויות $H$, $B$, $\mu_{\mathrm{r}}$ והיחס 
$r=\Delta V_x/\Delta V_y$
 מתקבלות משגיאה עקיפה לפי משוואה~\eqref{eq:HBpp} ו-$\mu_{\mathrm{r}}=B/(\mu_0 H)$.
% we should likely write them just in case.

\section{מעטפת שיא־לשיא של \texorpdfstring{$B(H)$}{B(H)}}

\subsection{מהלך הניסוי}
השנאי כוון 
ל-$\SI{10}{\volt},\SI{20}{\volt},\ldots,\SI{130}{\volt}$ ($13$ נקודות).
 בכל קפיצה הזזנו את הסמנים על האוסצילוסקופ עד שנגעו בקצות הלולאה, ורשמנו את הפרשי המתחים 
 $\Delta V_x,\Delta V_y$ (איור~\ref{fig:scope-cursors}).
  האמפרמטר שימש רק לבקרת מגבלת הזרם 
  ($\SI{2}{\ampere}$);
   הזרם בעיבוד חושב 
   מ-$I_x=V_x/R_x$.

\begin{figure}[H]
    \centering
    \includegraphics[width=0.82\linewidth]{hysteresis_schematic.pdf}
    \caption{המחשה לתהליך הניסוי}
    \label{fig:scope-cursors}
\end{figure}

\subsection{עיבוד נתונים}
כל זוג 
$(\Delta V_x,\Delta V_y)$
 הומר 
 ל-$(H,B)$
  דרך~\eqref{eq:HBpp}.
   התדריך מבקש עקומת מגנוט, אך בפועל נשמרו רק ממדי שיא־לשיא של הלולאות ולא ענף מגנוט ראשוני אחרי דה־מגנוט. לכן העיבוד כאן הוא תחליף ניסויי מוגבל: מעטפת שיאים שממנה ניתן לאמוד 
   $\mu_{\mathrm{r}}(H)$
    אפקטיבי ותחילת רוויה, אך לא רמננטיות, קוארסיביות או פרמאביליות ראשונית. הקווים החלקים שורטטו להנחיית העין בלבד, ו-
    $\mu_{\mathrm{r}}^{\max}$ נלקח כמקסימום בין נקודות המדידה. התוצאות באיור~\ref{fig:BH} ובטבלה~\ref{tab:bh-results}.

\begin{figure}[H]
    \centering
    \includegraphics[width=\linewidth]{fig_BH_mu.pdf}
    \caption{מעטפת השיאים
    $B(H)$
    והפרמאביליות היחסית האפקטיבית
    $\mu_{\mathrm{r}}(H)$
    שחושבה ממנה. הקווים החלקים הם להנחיית העין בלבד; הערכים הכמותיים מסוכמים בטבלה~%
    \ref{tab:bh-results}.}
    \label{fig:BH}
\end{figure}

\begin{table}[H]
\centering
\begin{tabular}{@{}rc@{}}
\toprule
פרמטר & ערך \\
\midrule
$\mu_{\mathrm{r}}^{\max}$ & $\num{2.18(8)e3}$ \\
$H(\mu_{\mathrm{r}}^{\max})$ & $\SI{167(4)}{\ampere\per\meter}$ \\
$B(\mu_{\mathrm{r}}^{\max})$ & $\SI{0.458(12)}{\tesla}$ \\
$B_{\mathrm{last}}$ & $\SI{1.243(33)}{\tesla}$ \\
\bottomrule
\end{tabular}
\caption{תוצאות מעטפת השיאים. הפרמאביליות היא אומדן מעטפת אפקטיבי, לא פרמאביליות ראשונית.}
\label{tab:bh-results}
\end{table}

העקומה עולה מהר בשדות קטנים ואז מתמתנת לקראת $B_{\mathrm{last}}\simeq\SI{1.24}{\tesla}$, כפי שמצופה כאשר מתחמים מגנטיים מתיישרים והחומר מתקרב לרוויה. $\mu_{\mathrm{r}}$ מגיע לשיא סביב $H\simeq\SI{167}{\ampere\per\meter}$ ואז יורד, כי תוספת שדה כבר מגדילה פחות את $B$. הערך $\mu_{\mathrm{r}}^{\max}\simeq\num{2.2e3}$ נמצא בתחום הגדלים הסביר לברזל רך, אך הוא כאמור ערך מעטפת אפקטיבי בלבד.

\section{פרמאביליות הריק}

\subsection{מהלך הניסוי}
בשתי סדרות $V_y$ ($\propto B$) הוחזק קבוע ולוחיות הנחושת הוסרו בזו אחר זו, כדי לשמור על $B$ קבוע בסריקת $L'$. ערך $B$ נבחר קרוב לרוויה תחת מגבלות הזרם. עוביי הרווח: $L'=\SIlist{0.35;0.20;0.10;0}{\milli\meter}$ בסדרה \texttt{a}, $\SIlist{0.50;0.45;0.35;0.20;0.10;0}{\milli\meter}$ בסדרה \texttt{b}. בכל סדרה נרשם אותו $\Delta V_y$ בכל נקודות הרווח, ולכן $B$ המחושב קבוע ברמת הדיוק שנרשמה (טבלה~\ref{tab:mu0-runs}). הערכים הקבועים $\langle B\rangle\simeq\SI{1.02}{\tesla}$ ו-$\SI{0.79}{\tesla}$ נמוכים מערך השיא של חלק א' ($\SI{1.24}{\tesla}$), ולכן אי־רוויה אפשרית נכללת במפורש כמקור שיטתי. פרמטרי הסדרות בטבלה~\ref{tab:mu0-runs}.

\begin{table}[H]
\centering
\begin{tabular}{@{}rcc@{}}
\toprule
פרמטר & סדרה \texttt{a} & סדרה \texttt{b} \\
\midrule
לוחיות התחלתיות & $3$ & $5$ \\
$\Delta V_y$ בכל נקודה & $\SI{3.67359}{\volt}$ & $\SI{2.83050}{\volt}$ \\
$\langle B\rangle$ & $\SI{1.0245}{\tesla}$ & $\SI{0.7894}{\tesla}$ \\
פיזור יחסי ב-$B$ & $<\num{0.01}\%$ & $<\num{0.01}\%$ \\
$H_{\max}$ & $\SI{1531}{\ampere\per\meter}$ & $\SI{1400}{\ampere\per\meter}$ \\
נקודות & $4$ & $6$ \\
\bottomrule
\end{tabular}
\caption{תנאי שתי סדרות הרווח במדידת
$\mu_0$.}
\label{tab:mu0-runs}
\end{table}

\subsection{עיבוד נתונים}
המודל~\eqref{eq:linear} לינארי $y=mx+b$ עם $y=NI/B$, $x=L'$, $m=1/\mu_0$, $b=L/\mu_{\mathrm{Fe}}$. עבור כל נקודה חושבו $I=\Delta V_x/(2R_x)$ ו-$B=R_yC\Delta V_y/(2NA)$, ואי־הוודאות של $y$ מהתפשטות שגיאות של $\Delta V_x,\Delta V_y,R_x,R_y,C,A$. אי־ודאות $L'$ מרזולוציית הקליבר. ערך \textenglish{CODATA}: $\mu_0^{\mathrm{theo}}=\SI{1.25664e-6}{\tesla\meter\per\ampere}$. ההתאמות באיור~\ref{fig:gap}, התוצאות בטבלה~\ref{tab:mu0-fit-results}.

\begin{figure}[H]
    \centering
    \includegraphics[width=0.9\linewidth]{fig_copper_gap.pdf}
    \caption{התאמת
    $\frac{NI}{B}$
    כפונקציה של
    $L'$
    בשתי סדרות המדידה. בחלק התחתון מוצג גרף השארים.}
    \label{fig:gap}
\end{figure}

\begin{table}[H]
\centering
\small
\begin{tabular}{@{}rcc@{}}
\toprule
פרמטר & סדרה \texttt{a} & סדרה \texttt{b} \\
\midrule
$m=\frac{1}{\mu_0}$ & $\SI{1.14(9)e6}{\ampere\per\tesla\per\meter}$ & $\SI{1.20(5)e6}{\ampere\per\tesla\per\meter}$ \\
$b=\frac{L}{\mu_{\mathrm{Fe}}}$ & $\SI{342(16)}{\ampere\per\tesla}$ & $\SI{220(12)}{\ampere\per\tesla}$ \\
$\mu_{0}^{\mathrm{fit}}=\frac{1}{m}$ & $\SI{8.8(7)e-7}{\tesla\meter\per\ampere}$ & $\SI{8.32(33)e-7}{\tesla\meter\per\ampere}$ \\
$\sigma_{\mu_0}/\mu_0$ & $8.2\%$ & $3.9\%$ \\
$\mu_{0}^{\mathrm{fit}}/\mu_{0}^{\mathrm{theo}}$ & $\num{0.70(6)}$ & $\num{0.66(3)}$ \\
$\mu_{\mathrm{Fe}}/\mu_{0}^{\mathrm{theo}}$ & $\num{1.12(5)e3}$ & $\num{1.74(10)e3}$ \\
\midrule
$\chi^2/\nu$ & $0.91$ & $0.67$ \\
$p$ & $0.403$ & $0.616$ \\
$\nu$ & $2$ & $4$ \\
$N_{\sigma}$ מול $\mu_{0}^{\mathrm{theo}}$ & $5.3$ & $13.0$ \\
\bottomrule
\end{tabular}
\caption{פרמטרי התאמת
$y=mx+b$
למדידת
$\mu_0$
וסטטיסטיקת ההתאמה.}
\label{tab:mu0-fit-results}
\end{table}

באיור~\ref{fig:gap} שתי הסדרות נראות לינאריות, וגם השארים אינם מצביעים על סטייה שיטתית בתוך כל סדרה. לכן הבעיה אינה איכות ההתאמה הלינארית: $\chi^2/\nu\sim1$ וערכי $p$ תקינים, אבל השיפועים נותנים $\mu_0$ נמוך מהספרות במובהקות גבוהה ($N_\sigma=5.3,13.0$). לכן נראה שמקור הפער בשגיאה שיטתית.

\section{דיון ומסקנות}
\begin{itemize}
\item מעטפת השיאים נתנה $\mu_{\mathrm{r}}^{\max}=\num{2.18(8)e3}$ — אומדן אפקטיבי, לא פרמאביליות ראשונית (טבלה~\ref{tab:bh-results}).
\item שתי סדרות פרמאביליות הריק לינאריות פנימית אך $\mu_0^{\mathrm{fit}}/\mu_0^{\mathrm{theo}}\approx0.7$ במובהקות גבוהה (טבלה~\ref{tab:mu0-fit-results}); ההסבר העיקרי הוא שטף/חתך אפקטיבי קטן מהגאומטרי, או רווח אפקטיבי גדול מהנומינלי, כך ש-$\alpha/\beta\approx0.7$ (טבלה~\ref{tab:mu0-systematics}).
\item במדידה חוזרת: למדוד עובי כל לוחית ישירות, לתעד סקאלות אוסילוסקופ, ולכייל עצמאית את האינטגרטור ואת השטח האפקטיבי.
\end{itemize}

\pagebreak{}

\part{שבוע 2}
\label{part:week2}

\section{מטרת הניסוי}
\begin{itemize}
\item מדידת מעבר פאזה פרומגנט-פאראמגנט בטמפרטורת קירי.
\end{itemize}

\section{רקע תיאורטי}

\subsection{מעבר פאזה}
מעל $T_{\mathrm{c}}$ תנודות תרמיות שוברות את הסידור הספונטני, והחומר עובר לפאזה פאראמגנטית ($M=0$ ב-$H=0$). מתחת ל-$T_{\mathrm{c}}$ קיימת מגנטיזציה ספונטנית. זהו מעבר רציף (סדר שני) שבו $M$ פרמטר הסדר.

\subsection{חוק קירי--וייס}
בפאזה הפאראמגנטית ($T>T_{\mathrm{c}}$),
\begin{equation}
\chi=\frac{M}{H}=\frac{C_{\mathrm{CW}}}{T-T_{\mathrm{c}}},
\end{equation}
ו-$1/\chi(T)$ לינארי, חותך את ציר $T$ ב-$T_{\mathrm{c}}$.
בפועל לא השתמשנו בהתאמת $1/\chi(T)$: הגאומטריה של הסלילים והדגימה לא נמדדה מספיק טוב להמרה אבסולוטית ל-$H,M$, ומעל המעבר נשארת תגובה מושרית בשדה סופי. לכן חוק קירי--וייס משמש כאן כרקע תאורטי בלבד, והחילוץ הכמותי נעשה ממדדי מגנוט מנורמלים.

\subsection{קירוב שדה ממוצע ליד המעבר}
מודל השדה הממוצע נותן
\begin{equation}
M(T)\propto\sqrt{T_{\mathrm{c}}-T}\qquad(T\lesssim T_{\mathrm{c}}),
\end{equation}
ולכן $M^2(T)$ לינארי סמוך ל-$T_{\mathrm{c}}$.

בשדה אפס פרמטר הסדר נעלם בחדות ב-$T_{\mathrm{c}}$; בשדה סופי המגנטיזציה המושרת לא מתאפסת בדיוק במעבר והמעבר מתרחב (איור~\ref{fig:curie-finite-field-schematic}). לכן הערך המחולץ הוא טמפרטורת חצי־גובה תפעולית בשדה סופי, ולא $T_{\mathrm{c}}$ תרמודינמי בשדה אפס; סדרות בשדות שונים אינן חזרות שקולות.

\begin{figure}[H]
    \centering
    \includegraphics[width=0.58\linewidth]{curie_finite_field_schematic.pdf}
    \caption{סכמה של דעיכת המגנוט ליד מעבר קירי בשדה אפס ובשדה חיצוני סופי. שדה גדול יותר מרחיב את המעבר ומותיר מגנטיזציה מושרית גם מעל
    $T_{\mathrm{c}}$.}
    \label{fig:curie-finite-field-schematic}
\end{figure}

\section{מערכת הניסוי}
המערכת דומה לשבוע~1: סליל ראשוני נותן $H$ וסליל משני מחובר לאינטגרטור $RC$ עם $V_y\propto B$. הדגימה ממוקמת במכל קירור חנקן נוזלי, וטמפרטורתה נמדדת בצמד־חום המחובר ל-\textenglish{LabVIEW}. הדגימה מקוררת ואז מתחממת בהדרגה לטמפרטורת החדר; $X,Y,T$ נדגמים ברציפות (איור~\ref{fig:curie-circuit}).

\begin{figure}[H]
    \centering
    \includegraphics[width=0.8\linewidth]{experiment2.jpeg}
    \caption{מערכת ניסוי טמפרטורת קירי: דגימת מונל מקוררת בחנקן נוזלי ונמדדת בעזרת סלילי עירור ומדידה.}
    \label{fig:curie-circuit}
\end{figure}

\subsection{ציוד}
\begin{itemize}
\item דגימת מונל בעלת מעבר מגנטי בתחום הנמדד.
\item מכל חנקן נוזלי מבודד תרמית.
\item צמד־חום ומד טמפרטורה \textenglish{508BR} מחובר ל-\textenglish{LabVIEW}.
\item שנאי משתנה, נגד $R_x$, אינטגרטור $RC$ — כבשבוע~1.
\item אוסילוסקופ \textenglish{Agilent InfiniiVision DSO7012A} ומולטימטר \textenglish{HP 34401A}.
\end{itemize}

רזולוציות הצירים בטבלאות~\ref{tab:curie-voltage-sigmas} ו-\ref{tab:curie-sigmaT}.

\subsection{סימוני הצירים והכיול היחסי}
המעגל זהה לשבוע~1, עם קבועים $N_1=250$, $N_2=2500$, $R_y=\SI{3.97}{\kilo\ohm}$, $C=\SI{19.78}{\microfarad}$. אורך הסולנואיד ושטח החתך לא נמדדו ישירות, ולכן לא ניתן לבצע המרה אבסולוטית ל-$H,M$. לכן לא מדווחים כאן מגנטיזציה אבסולוטית, אלא עוקבים אחרי מדד מגנוט יחסי כפונקציה של טמפרטורה. נשמרים דגמי הענפים עצמם, $X\equiv V_x\propto H$, $Y\equiv V_y\propto B$.

הזרמים: $I_{\mathrm{rms}}=\SI{1.97}{\ampere}$ ב-\textenglish{A}, $\SI{2.24}{\ampere}$ ב-\textenglish{B}, $\SI{0.65}{\ampere}$ ב-\textenglish{C}. הקואורדינטות בעיבוד יחסיות,
\begin{equation}
H^*=c_H X,\qquad M^*=c_Y Y-H^*,
\end{equation}
עם $c_H,c_Y$ קבועי סקאלה לכל סדרה. איורי הנתונים הגולמיים בוולטים, ואחרי ההמרה ב-$H^*,M^*$. חילוץ $T_{1/2}^{\mathrm{app}}$ נשען על הדעיכה המנורמלת בלבד.

\section{שיטת מדידה ועיבוד נתונים}

\subsection{מהלך הניסוי}
הדגימה קוררה לחנקן נוזלי ואז הושארה להתחמם בהדרגה לטמפרטורת החדר. $X,Y,T$ נדגמו ברציפות; כל מחזור $50\,\si{\hertz}$ נשמר כלולאה. סדרה~\textenglish{A} שימשה כסדרה הראשית, \textenglish{B} נמדדה בשדה גבוה יותר, ו-\textenglish{C} כבדיקת רגישות בזרם נמוך.

הלולאות הגולמיות מראות את המעבר: בטמפרטורה נמוכה לולאת חשל פתוחה, סוגרת עם החימום, ובסביבת $\SI{270}{\kelvin}$ נשארת תגובה כמעט לינארית (איור~\ref{fig:curie-scope-transition}).

\begin{figure}[H]
    \centering
    \includegraphics[width=0.92\linewidth]{curie_scope_transition_frames.pdf}
    \caption{לולאות מייצגות מסדרה~%
    \textenglish{A}
    באותם צירים, לאורך החימום. התפתחות הצורה מלולאה פתוחה לתגובה כמעט לינארית מראה את דעיכת המגנוט; לא כל הלולאות המוצגות משמשות בהמשך לעיבוד הכמותי.}
    \label{fig:curie-scope-transition}
\end{figure}

סדר העיבוד היה:
\begin{itemize}
\item לנקות את הסדרות מלולאות לא יציבות בתחילת המדידה ומעודף נקודות חמות בסוף.
\item להוציא מכל לולאה שלושה מדדי מגנוט יחסיים: רמננטיות, קצה רוויה, ואקסטרפולציית זנבות.
\item לנרמל כל מדד ולקרוא את טמפרטורת חצי־הגובה.
\item להשוות בין הסדרות כדי להעריך רגישות לשדה ולכיסוי הטמפרטורה.
\end{itemize}

\subsection{הכנת הנתונים}
הסינון נעשה בשני שלבים, כדי שלא לערבב שינוי פיזיקלי במגנוט עם שינוי טכני של המדידה:

\textbf{שלב~1: חיתוך הקצה החם.} שלוש הסדרות קוצרו לסיום משותף ב-$T\simeq\SI{273.3}{\kelvin}$, כדי למנוע ייצוג עודף של נקודות חמות בחיסור הרקע ובנרמול.

\textbf{שלב~2: חיתוך התייצבות השדה.} בלולאות הראשונות אמפליטודת השדה עדיין מתייצבת, ולכן קצות הלולאה אינם בני־השוואה לשאר הסדרה. כדי להחליט מאיזו לולאה להתחיל, הגדרנו חלון שדה משותף לענפים $(\pm)$:
\begin{equation}
X_{\mathrm{win}}=\min\!\left(\max X^{(+)},\,\bigl|\min X^{(-)}\bigr|\right),
\end{equation}
ו-$X_{\mathrm{plateau}}$ חציון $X_{\mathrm{win}}$ במחצית השנייה של הסדרה. נשמרה הלולאה הראשונה שבה $X_{\mathrm{win}}\ge0.98\,X_{\mathrm{plateau}}$. דוגמה ללולאות שנפסלו באיור~\ref{fig:curie-rejected-loop-example}.

לאחר השלבים: \textenglish{A} עם $141$ לולאות ב-$\SI{180.8}{\kelvin}$--$\SI{273.3}{\kelvin}$; \textenglish{B} ב-$\SI{159.1}{\kelvin}$--$\SI{273.2}{\kelvin}$; \textenglish{C} ב-$\SI{221.3}{\kelvin}$--$\SI{273.3}{\kelvin}$.

\begin{figure}[H]
    \centering
    \includegraphics[width=0.62\linewidth]{curie_rejected_loop_example.pdf}
    \caption{דוגמה מסדרה~%
    \textenglish{A}
    ללולאות שנפסלו בחיתוך התייצבות השדה, לעומת הלולאה הראשונה שנשמרה. בלולאות שנפסלו חלון השדה וקצות הלולאה עדיין משתנים, וזנבות רוויה יציבים אינם מוגדרים היטב; לכן הן אינן מתאימות לחילוץ מדד הרוויה או להתאמת זנבות הרוויה של שיטה~3.}
    \label{fig:curie-rejected-loop-example}
\end{figure}

איור~\ref{fig:curie-data-prep-cuts} מסכם את החיתוך בשלוש הסדרות. סדרה~\textenglish{C}, לאחר החיתוך, חסרה את התחום הקר היציב מתחת למעבר, ולכן נשמרת כבדיקת רגישות לשדה ולכיסוי טמפרטורה ולא כחזרה שקולה.

\begin{figure}[H]
    \centering
    \includegraphics[width=0.86\linewidth]{curie_data_preparation_cuts.pdf}
    \caption{חיתוך התייצבות השדה בשלוש סדרות קירי לפי הקריטריון
    $X_{\mathrm{win}}/X_{\mathrm{plateau}}\ge0.98$.
    הנקודות האפורות נפסלו, והנקודות הצבעוניות נשמרו לעיבוד.}
    \label{fig:curie-data-prep-cuts}
\end{figure}

\subsection{חילוץ \texorpdfstring{$M(T)$}{M(T)} מן הלולאות}
לפני החילוץ הוסר רקע לינארי שנקבע מן הרבעון העליון של טמפרטורות הסדרה (רקע אמפירי, לא הנחת תחום קירי--וייס נקי). $M$ או $M^*$ להלן הם מדד מגנוט יחסי לאחר חיסור הרקע.

השתמשנו בשלושה מדדים מאותה לולאה:
\begin{itemize}
\item $M_{\mathrm{r}}$: פתיחת הלולאה ב-$H^*=0$.
\item $M_{\mathrm{sat}}$: פתיחת הלולאה בקצה חלון השדה המשותף.
\item $M_0$: אקסטרפולציה של זנבות הרוויה אל $H^*=0$.
\end{itemize}

בכל טמפרטורה ללולאה שני ענפים. הסימונים הבאים אינם נקודת ענף יחידה אלא חצי הפרש סימטרי בין הענפים, מה שמבטל היסט משותף בציר האנכי:
\begin{align}
M_{\mathrm{r}}(T)&=\frac{1}{2}\left|M^*_{+}(T;0)-M^*_{-}(T;0)\right|,\\
M_{\mathrm{sat}}(T)&=\frac{1}{2}\left|M^*_{+}(T;H^*_{\mathrm{sat}})-M^*_{-}(T;-H^*_{\mathrm{sat}})\right|,\\
M_{0}(T)&=\frac{1}{2}\left|b^{(+)}-b^{(-)}\right|.
\end{align}
בשיטה~3, $b^{(\pm)}$ הם חיתוכי ההתאמה הלינארית בזנבות הרוויה, $M^*(H^*)\approx M_{0}(T)+\chi_{\mathrm{HF,rel}}(T)H^*$ — הצורה היחסית של $B/\mu_{0}-(\chi_{\mathrm{HF}}+1)H=M_{0}$. חלון ההתאמה נבחר מזנבות הלולאה בלבד: מתחילים מחמש נקודות הקצה, שומרים נקודות עם $|M^*|\ge0.95|\hat M_0|$, ומרחיבים פנימה כל עוד השיפוע נשאר יציב.

כל שלושת הגדלים נורמלו באותה צורה:
\begin{equation}
y_{\mathrm{norm}}(T)=\frac{y(T)-\min_T y}{\max_T y-\min_T y}.
\end{equation}

הדגמה של השיטות על לולאה אחת באיור~\ref{fig:curie-method3-tail-example}; המדדים המנורמלים באיור~\ref{fig:curie-method123}.

\begin{figure}[H]
    \centering
    \includegraphics[width=0.94\linewidth]{curie_method3_tail_example.pdf}
    \caption{הדגמה גאומטרית של שלוש שיטות החילוץ על לולאת חשל אחת, במתחי האוסילוסקופ
    \textenglish{\(V_x,V_y\)}:
    שיטה~1 קוראת את פתיחת הלולאה ב-%
    \textenglish{\(V_x=0\)},
    שיטה~2 קוראת את קצות הרוויה ב-%
    \textenglish{\(V_x=\pm V_{x,\mathrm{sat}}\)},
    ושיטה~3 מתאימה את זנבות הרוויה וממשיכה את הקווים ל-%
    \textenglish{\(V_x=0\)}.
    בעיבוד המספרי אותן קריאות מומרות לינארית ל-%
    \textenglish{\(H^*,M^*_{\mathrm{corr}}\)};
    חצאי ההפרשים נותנים את
    \textenglish{\(M_r,M_{\mathrm{sat}},M_0\)}.}
    \label{fig:curie-method3-tail-example}
\end{figure}

\begin{figure}[H]
    \centering
    \includegraphics[width=0.85\linewidth]{curie_method123_normalized.pdf}
    \caption{שלושת מדדי הסדר המנורמלים בסדרה~%
    \textenglish{A}.
    הקו האופקי מסמן חצי־גובה, והסימנים עליו מציינים את חציית כל מדד.}
\label{fig:curie-method123}
\end{figure}

באיור~\ref{fig:curie-method123} שתי השיטות הישירות יורדות כמעט באותו תחום טמפרטורות, ולכן הן בסיס יציב יותר לאומדן המרכזי. שיטה~3 יורדת מוקדם יותר, מפני שהיא רגישה לבחירת זנבות הרוויה ולחיסור הרקע. לכן היא נשמרת בבדיקת אי־הוודאות, אך לא מתפרשת כחזרה בלתי־תלויה מלאה.

\subsection{אי־ודאויות}
בנרמול לחצי־גובה, כל סקאלה משותפת של $M$ ($N,L,R_x,R_y,C,A$) מבטלת את עצמה. התקציב מתחלק לשני סוגים: אי־ודאות טמפרטורה מוחלטת ללולאה (טבלה~\ref{tab:curie-sigmaT}) ואי־ודאות חילוץ מקומית מצירי המתח ומההתאמות (טבלה~\ref{tab:curie-voltage-sigmas}). האי־ודאויות בסוגריים בטבלאות התוצאות הן מקומיות בלבד; דיוק מד הטמפרטורה המוחלט הוא תרומה שיטתית משותפת ונבלע במעטפת הסופית.

תקציב הטמפרטורה הוא
\begin{equation}\label{eq:sigmaT-curie}
\sigma_T=\sqrt{\sigma_{T,\mathrm{smear}}^2+\sigma_{T,\mathrm{res}}^2+\sigma_{T,\mathrm{instr}}^2(T)}.
\end{equation}
ב-$T_{1/2}^{\mathrm{app}}\approx\SI{215}{\kelvin}$ ($\SI{-58}{\degreeCelsius}$) תקף התחום החם של המפרט, $\sigma_{T,\mathrm{instr}}\approx\SI{1.1}{\kelvin}$; מתחת ל-$\SI{-60}{\degreeCelsius}$, $\sigma_{T,\mathrm{instr}}\approx\SI{2.1}{\kelvin}$. המפרט מתייחס למד הטמפרטורה בלבד ואינו כולל את צמד־החום, ולכן הוא אינו מצדיק דיוק מוחלט טוב יותר מ-$\sim\SI{1}{\kelvin}$--$\SI{2}{\kelvin}$.

\begin{table}[H]
\centering
\small
\setlength{\tabcolsep}{3pt}
\begin{tabular}{@{}r >{\centering\arraybackslash}p{0.32\linewidth} >{\centering\arraybackslash}p{0.43\linewidth}@{}}
\toprule
תרומה & מקור & הגדרה \\
\midrule
$\sigma_{T,\mathrm{smear}}$ & טשטוש חימום בחלון $\Delta t_{\mathrm{loop}}\approx\SI{5.6}{\second}$ & $\displaystyle\frac{|dT/dt|\,\Delta t_{\mathrm{loop}}}{\sqrt{12}}$ \\
$\sigma_{T,\mathrm{res}}$ &  שגיאת רזולוציה & $\displaystyle\frac{\SI{1e-3}{\kelvin}}{\sqrt{12}}$ \\
$\sigma_{T,\mathrm{instr}}$ & שגיאת מכשיר & $\displaystyle
\begin{cases}
0.1\%\,|T_{^\circ\mathrm{C}}|+\SI{1}{\kelvin}, & T\in[\SI{-60}{\degreeCelsius},\SI{1372}{\degreeCelsius}],\\
0.1\%\,|T_{^\circ\mathrm{C}}|+\SI{2}{\kelvin}, & T\in[\SI{-220}{\degreeCelsius},\SI{-60}{\degreeCelsius}).
\end{cases}$ \\
\bottomrule
\end{tabular}
\caption{אי־ודאות הטמפרטורה לכל לולאה.}
\label{tab:curie-sigmaT}
\end{table}

לכל שיטה הוערכה אי־ודאות חילוץ מקומית $\sigma_y$ מן ההתאמה הלינארית (סביב $X=0$ בשיטה~1, סביב $X_{\min/\max}$ בשיטה~2, בזנב הרוויה בשיטה~3), בעזרת אי־ודאויות הדיגיטציה (טבלה~\ref{tab:curie-voltage-sigmas}). ההתאמה מתבצעת במתחים, וההמרה ל-$M^*(H^*)$ רק בסיום.

\begin{table}[H]
\centering
\small
\begin{tabular}{@{}rcc@{}}
\toprule
תרומה & מקור & גודל \\
\midrule
$\sigma_{X,\mathrm{dig}}$ & שגיאת דיגיטציה בציר $X$, $q_x=\SI{0.6}{\milli\volt}$ & $q_x/\sqrt{12}$ \\
$\sigma_{Y,\mathrm{dig}}$ & שגיאת דיגיטציה בציר $Y$, $q_y=\SI{0.4}{\milli\volt}$ & $q_y/\sqrt{12}$ \\
\bottomrule
\end{tabular}
\caption{אי־ודאות צירי המתח בהתאמות זנבות הלולאה.}
\label{tab:curie-voltage-sigmas}
\end{table}

\section{תוצאות טמפרטורת חצי־גובה}

\subsection{חילוץ טמפרטורת חצי־הגובה}
לכל מדד חושבה הטמפרטורה שבה המדד המנורמל יורד דרך חצי־גובה. אי־הוודאות המדווחת בטבלאות היא אי־ודאות חילוץ מקומית של החציה, ואינה כוללת את דיוק מד הטמפרטורה המוחלט. תוצאות סדרה~\textenglish{A} (טבלה~\ref{tab:curie-results}) קובעות את האומדן המרכזי ואת רגישות השיטה; טבלה~\ref{tab:curie-runs} משמשת לבדיקת רגישות לשדה ולכיסוי טמפרטורה.

\begin{table}[H]
\centering
\begin{tabular}{@{}rcc@{}}
\toprule
מדד & $T_{1/2}^{\mathrm{app}}$ (\si{\kelvin}) & פירוש \\
\midrule
$M_{\mathrm{r}}$ & $\num{216.39(77)}$ & רמננטיות \\
$M_{\mathrm{sat}}$ & $\num{214.76(95)}$ & קצה השדה המשותף \\
$M_{0}$ & $\num{199.14(169)}$ & אקסטרפולציה \\
\bottomrule
\end{tabular}
\caption{חציות חצי־הגובה בסדרה~%
\textenglish{A};
אי־הוודאות מקומית בלבד.}
\label{tab:curie-results}
\end{table}

שתי השיטות הישירות בסדרה~\textenglish{A}, $M_{\mathrm{r}}$ ו-$M_{\mathrm{sat}}$, מסכימות בתוך כ-$\SI{2}{\kelvin}$. $M_0$ נמוך יותר בכ-$\SI{16}{\kelvin}$, ולכן הוא מסמן רגישות לשיטת אקסטרפולציה ולא חזרה שקולה על אותה מדידה.

שלוש השיטות חושבו מחדש על שתי הסדרות הנוספות. טבלה~\ref{tab:curie-runs} משווה בין הסדרות והשיטות. שינוי הנגד הראשוני משנה את אמפליטודת השדה ואת צורת הלולאה, ולכן הסדרות אינן חזרות שקולות אלא בדיקת רגישות לשדה ולכיסוי הטמפרטורה.

\begin{table}[H]
\centering
\footnotesize
\setlength{\tabcolsep}{3pt}
\newcommand{\curieresult}[2]{\ensuremath{#1\pm #2}}
\newcommand{\relerr}[1]{\ensuremath{#1\%}}
\begin{tabular}{@{}rccccccc@{}}
\toprule
סדרה & $I_{\mathrm{rms}}$ (\si{\ampere}) & $H/H_A$ & תחום $T$ (\si{\kelvin}) & לולאות & מדד & $T_{1/2}\pm\sigma$ (\si{\kelvin}) & $\sigma/T$ \\
\midrule
\textenglish{A} & $1.97$ & $1.00$ & $180.8$--$273.3$ & $141$ & $M_{\mathrm{r}}$ & \curieresult{216.4}{0.8} & \relerr{0.4} \\
\textenglish{A} & $1.97$ & $1.00$ & $180.8$--$273.3$ & $141$ & $M_{\mathrm{sat}}$ & \curieresult{214.8}{1.0} & \relerr{0.5} \\
\textenglish{A} & $1.97$ & $1.00$ & $180.8$--$273.3$ & $141$ & $M_0$ & \curieresult{199.1}{1.7} & \relerr{0.9} \\
\addlinespace[0.25em]
\textenglish{B} & $2.24$ & $1.13$ & $159.1$--$273.2$ & $155$ & $M_{\mathrm{r}}$ & \curieresult{209.7}{2.1} & \relerr{1.0} \\
\textenglish{B} & $2.24$ & $1.13$ & $159.1$--$273.2$ & $155$ & $M_{\mathrm{sat}}$ & \curieresult{210.0}{1.7} & \relerr{0.8} \\
\textenglish{B} & $2.24$ & $1.13$ & $159.1$--$273.2$ & $155$ & $M_0$ & \curieresult{192.5}{1.6} & \relerr{0.8} \\
\addlinespace[0.25em]
\textenglish{C} & $0.65$ & $0.33$ & $221.3$--$273.3$ & $177$ & $M_{\mathrm{r}}$ & \curieresult{245.3}{0.8} & \relerr{0.3} \\
\textenglish{C} & $0.65$ & $0.33$ & $221.3$--$273.3$ & $177$ & $M_{\mathrm{sat}}$ & \curieresult{233.6}{0.8} & \relerr{0.3} \\
\textenglish{C} & $0.65$ & $0.33$ & $221.3$--$273.3$ & $177$ & $M_0$ & \curieresult{229.4}{0.8} & \relerr{0.3} \\
\bottomrule
\end{tabular}
\caption{סדרות קירי מול שלוש שיטות חצי־גובה, עם אי־ודאות חילוץ מקומית.
כל שורה היא סדרה ושיטת חילוץ אחת; שתי העמודות האחרונות מציגות את אי־הוודאות האבסולוטית ואת השגיאה היחסית $100\sigma/T$.
$H/H_A$
הוא יחס אמפליטודת השדה החציונית לסדרה~%
\textenglish{A};
סדרה~%
\textenglish{C}
מדווחת כמעטפת רגישות.}
\label{tab:curie-runs}
\end{table}

השיטות הן בחירות עיבוד שונות ולא חזרות שקולות של אותה מדידה. בשתי השיטות הישירות סדרה~\textenglish{A} נותנת $\SI{216.4}{\kelvin}$ ו-$\SI{214.8}{\kelvin}$, וסדרה~\textenglish{B} נותנת $\SI{209.7}{\kelvin}$ ו-$\SI{210.0}{\kelvin}$. סדרה~\textenglish{C} (שדה $0.33H_A$, ללא תחום קר יציב) נותנת ערכים גבוהים יותר, ונשמרת כמעטפת רגישות רחבה.

איור~\ref{fig:curie-drive-comparison} מציג את $M^*_{\mathrm{sat}}$ המנורמל בשלוש הסדרות מול ציר טמפרטורה מנורמל. שלוש העקומות אינן מתלכדות לעקומה תאורטית יחידה — סימן שגם השהיה תרמית, תחום רוויה, חזרתיות וכיסוי טמפרטורה משפיעים על התוצאה.

\begin{figure}[H]
    \centering
    \includegraphics[width=0.82\linewidth]{curie_measured_drive_comparison.pdf}
    \caption{דעיכת מדד הרוויה היחסי
    $M^*_{\mathrm{sat}}$
    המנורמל בשלוש סדרות המדידה. ערכי הזרם והכיסוי של כל סדרה מסוכמים בטבלה~%
    \ref{tab:curie-runs}.}
    \label{fig:curie-drive-comparison}
\end{figure}

\section{דיון ומסקנות}
הרכב סגסוגת המונל אינו ידוע, ולכן איכות התוצאה נבחנת מעקביות פנימית: הסכמה בין שלוש שיטות, חזרתיות בין סדרות, ורגישות לשדה ולכיסוי.

סדרה~\textenglish{A} היא הבסיס: היא כוללת תחום קר יציב, כיסוי מלא עד הפאראמגנטי, ושתי השיטות הישירות ($M_{\mathrm{r}},M_{\mathrm{sat}}$) מסכימות בתוך $\SI{2}{\kelvin}$. שיטה~3 ($M_0$) רגישה לבחירת חלון זנב הרוויה ולכן נותנת ערך נמוך יותר; הפיזור בינה לשתי הישירות נכלל באי־הוודאות. סדרה~\textenglish{B} נמוכה בכ-$\SI{6}{\kelvin}$ בשיטות הישירות ומשמשת לאומדן רגישות שדה בין סדרות שדה גבוה. סדרה~\textenglish{C} אינה שקולה (שדה קטן פי שלושה, חסר תחום קר) ונשמרת כמעטפת רחבה.

הערך המרכזי נלקח משתי השיטות הישירות בסדרה~\textenglish{A}. אי־הוודאות הסופית משלבת את פיזור השיטות בסדרה~\textenglish{A} עם פיזור הסדרות \textenglish{A,B} בשיטות הישירות:
\begin{equation}\label{eq:T12-headline}
T_{1/2}^{\mathrm{app}}=\SI{216(10)}{\kelvin}=\SI{-58(10)}{\degreeCelsius},
\end{equation}
כאשר $\sigma_{\mathrm{final}}=\sqrt{(\SI{9.5}{\kelvin})^2+(\SI{4.0}{\kelvin})^2}\simeq\SI{10}{\kelvin}$. עם הכללת \textenglish{C} מתקבלת מעטפת רחבה $\SI{216(20)}{\kelvin}$. בשני המקרים זהו מעבר תפעולי בשדה סופי, לא $T_{\mathrm{c}}$ תרמודינמי.

\begin{itemize}
\item שלושת המדדים בסדרה~\textenglish{A} יורדים באותו תחום; שיטה~3 רגישה יותר לחלון זנב הרוויה.
\item מקורות הטיה: השהיה וגרדיאנט תרמיים, חיסור רקע לינארי, חיתוך התייצבות, ובחירת חלון זנב.
\item במדידה חוזרת: למדוד את גאומטריית הדגימה, לשפר הצמדה תרמית, לסרוק לאט יותר, לכסות תחום קר יציב בכל סדרה, ולכייל את ערוץ $B$ עצמאית.
\end{itemize}

\section{סיכום כולל}
\begin{table}[H]
\centering
\small
\setlength{\tabcolsep}{3pt}
\begin{tabular}{@{}rccc@{}}
\toprule
תוצאה & ערך & מגבלה עיקרית & שיפור במדידה חוזרת \\
\midrule
$\mu_{\mathrm{r}}^{\max}$ & $\num{2.18(8)e3}$ & מעטפת שיאים, לא ענף מגנוט ראשוני & ייצוא לולאה מלאה ודה־מגנוט לפני סריקה \\
$\mu_0^{\mathrm{fit}}/\mu_0^{\mathrm{theo}}$ & $\num{0.70(6)},\num{0.66(3)}$ & שטף/חתך אפקטיבי ורווח נחושת & כיול $A_{\mathrm{eff}}$, מדידת כל לוחית וכיול ערוץ $B$ \\
$T_{1/2}^{\mathrm{app}}$ & $\SI{216(10)}{\kelvin}$ & שדה סופי, השהיה תרמית ובחירת שיטה & סריקה איטית, הצמדה תרמית טובה וכיסוי קר בכל סדרה \\
\bottomrule
\end{tabular}
\caption{סיכום תוצאות שני השבועות והגורמים המגבילים העיקריים.}
\label{tab:global-summary}
\end{table}

\end{document}
